\documentclass[11pt,a4paper]{article}

\usepackage[utf8]{inputenc}
\usepackage[T1]{fontenc}
\usepackage{lmodern}
\usepackage[margin=0.82in]{geometry}
\usepackage[table]{xcolor}
\usepackage{array}
\usepackage{tabularx}
\usepackage{booktabs}
\usepackage{enumitem}
\usepackage{titlesec}
\usepackage{fancyhdr}
\usepackage{lastpage}
\usepackage{hyperref}
\usepackage{tikz}

\definecolor{TextInk}{HTML}{233142}
\definecolor{HeadingBlue}{HTML}{1D3557}
\definecolor{AccentTeal}{HTML}{2A7F92}
\definecolor{AccentGold}{HTML}{C58B4E}
\definecolor{RuleGray}{HTML}{D6E0E5}
\definecolor{TableStripe}{HTML}{F3F8FA}
\definecolor{SoftTint}{HTML}{F7F4EE}
\definecolor{SoftBlue}{HTML}{EEF6F8}
\definecolor{SoftGreen}{HTML}{EEF7F0}
\definecolor{SoftRed}{HTML}{FBF0EF}

\hypersetup{
  colorlinks=true,
  linkcolor=HeadingBlue,
  urlcolor=AccentTeal,
  pdftitle={IB Geography SL Study Pack - 2026-04-29},
  pdfauthor={IB Geography SL Preparation}
}

\color{TextInk}
\setlength{\parskip}{0.52em}
\setlength{\parindent}{0pt}
\setlength{\tabcolsep}{6pt}
\renewcommand{\arraystretch}{1.22}
\setlist[itemize]{leftmargin=1.35em,itemsep=3pt,topsep=4pt}
\setlist[enumerate]{leftmargin=1.5em,itemsep=3pt,topsep=4pt}

\newcolumntype{Y}{>{\raggedright\arraybackslash}X}
\newcolumntype{L}[1]{>{\raggedright\arraybackslash}p{#1}}

\titleformat{\section}
  {\normalfont\huge\sffamily\bfseries\color{HeadingBlue}}
  {}{0pt}{}
  [\vspace{0.15em}{\color{AccentGold}\titlerule[1pt]}]

\titleformat{\subsection}
  {\normalfont\Large\sffamily\bfseries\color{AccentTeal}}
  {}{0pt}{}

\titleformat{\subsubsection}
  {\normalfont\large\sffamily\bfseries\color{HeadingBlue}}
  {}{0pt}{}

\titlespacing*{\section}{0pt}{1.4em}{0.65em}
\titlespacing*{\subsection}{0pt}{1.0em}{0.35em}
\titlespacing*{\subsubsection}{0pt}{0.8em}{0.25em}

\pagestyle{fancy}
\setlength{\headheight}{14pt}
\fancyhf{}
\fancyhead[L]{\sffamily\color{AccentTeal}Wed 2026-04-29}
\fancyhead[R]{\sffamily\color{HeadingBlue}IB Geography SL}
\fancyfoot[C]{\sffamily\color{HeadingBlue}\thepage\ of \pageref{LastPage}}
\renewcommand{\headrulewidth}{0.4pt}
\renewcommand{\footrulewidth}{0pt}

\newcommand{\term}[1]{\textcolor{HeadingBlue}{\textbf{#1}}}
\newcommand{\exam}[1]{\textcolor{AccentTeal}{\textbf{#1}}}
\newcommand{\boxnote}[3]{%
\noindent\colorbox{#1}{%
  \begin{minipage}{0.965\textwidth}
  \sffamily
  \textbf{\textcolor{HeadingBlue}{#2}} #3
  \end{minipage}
}}

\begin{document}

\begin{center}
  {\sffamily\bfseries\color{HeadingBlue}\LARGE IB Geography SL Study Pack}\\[0.25em]
  {\sffamily\bfseries\color{AccentTeal}\Large Wednesday 2026-04-29}\\[0.35em]
  {\sffamily\color{TextInk}Urban Environments Essay Work + Population Distribution And Changing Population}
\end{center}

\vspace{0.5em}

\boxnote{SoftBlue}{How this pack follows the plan.}{The April 29 schedule is a
weekday session. \term{Audio 1} covers \term{Urban environments} essay planning
strategies, social and environmental stresses, sustainability, and future urban
management debate. \term{Audio 2} revisits \term{Population distribution --
changing population}. The \term{25 minute} practice block asks you to plan one
\term{Paper 1 extended answer} from a \term{choice of 2} prompts, with a clear
judgment. Use the separate April 29 essay-planning workbook for the timed plan.}

\section{Today's Target}

By the end of this session, you should be able to:

\begin{itemize}
  \item turn an Urban environments extended-answer prompt into a focused line of argument
  \item separate \term{social stresses}, \term{environmental stresses}, and
  \term{management responses} before writing
  \item evaluate future urban management using stakeholders, scale, trade-offs,
  and sustainability
  \item use \term{Detroit} carefully for gentrification, deindustrialization,
  selective reinvestment, and social inequality
  \item use a second urban management example from class notes, or a clear
  placeholder if the exact case still needs to be filled
  \item recall the population frame: \term{pattern, process, structure, policy}
  \item explain population change using births, deaths, and migration
  \item compare population policy examples such as China, Japan, and Singapore
  \item plan a \term{10-mark} style answer with an explicit judgment before writing
\end{itemize}

\subsection{Session Structure}

\rowcolors{2}{TableStripe}{white}
\begin{tabularx}{\textwidth}{L{0.18\textwidth}Y L{0.25\textwidth}}
\toprule
\textbf{Part} & \textbf{Focus} & \textbf{Output} \\
\midrule
Audio 1 & Urban essay planning, stresses, sustainability, future city debates &
One argument map for Urban environments \\
Audio 2 & Population distribution, change, structure, policy, and examples &
One population comparison chain \\
25 min practice & Choose one of two Paper 1 extended-answer prompts and plan it &
One timed essay plan with a clear judgment \\
\bottomrule
\end{tabularx}
\rowcolors{2}{}{}

\boxnote{SoftTint}{Carry-forward from April 28.}{Bring forward yesterday's
visual-data weakness. If the issue was weak evidence, put case-study details
into the plan. If the issue was description without explanation, make every
urban paragraph use the chain \term{stress} $\rightarrow$ \term{cause}
$\rightarrow$ \term{impact} $\rightarrow$ \term{response} $\rightarrow$
\term{evaluation}.}

\section{Audio 1 Script}

\subsection{Urban Environments: Essay Planning, Stresses, Sustainability, And Future Management}

\textit{Target length: about 12 minutes at a natural revision pace. Pause at
each planning checkpoint and say the answer structure aloud before moving on.}

\textbf{Minute 0--2: what an extended answer needs.}
Today the Urban environments target is essay work. That means the job is not to
collect more facts. The job is to turn facts into an argument. In Paper 1, the
extended answer is where you show selection, structure, evaluation, and
judgment. You still need accurate content, but content by itself is not enough.
The examiner needs to see that you know which evidence matters for the exact
question.

Start every extended answer plan with three questions. First, what is the
command term? \exam{Evaluate} asks for strengths, weaknesses, and a supported
judgment. \exam{Examine} asks you to look closely at more than one side of an
issue. \exam{To what extent} asks how far a statement is true, so your answer
must contain a degree of agreement. Second, what is the topic limit? If the
question says \term{social and environmental stresses}, do not write a general
history of urbanization. If it says \term{sustainable urban management}, do not
only describe problems. Third, what is your judgment? A judgment is not just a
final sentence saying "overall." It is the line of argument running through the
whole answer.

Use this planning sentence: \term{Although} one side is important, \term{the
stronger judgment is} ... \term{because} ... For example: "Although sustainable
urban management can reduce some environmental stresses, its success is uneven
because social inequality, land value, and governance decide who benefits."
That sentence already contains balance, judgment, and criteria.

\textbf{Memory check.} Name three command terms that need evaluation. Then say
one planning sentence beginning with "Although".

\textbf{Planning rule before content.}
Before you choose examples, build the answer frame. A reliable extended-answer
paragraph has five parts: \term{point}, \term{evidence}, \term{explanation},
\term{evaluation}, and a \term{mini-judgment}. The point answers part of the
question. The evidence names a place, strategy, or process. The explanation
shows how it works. The evaluation says why it is successful, limited, uneven,
or contested. The mini-judgment links the paragraph back to the line of
argument. If a paragraph has only point and evidence, it becomes description. If
it has point, evidence, explanation, and evaluation, it becomes an argument.

Here is the difference. A descriptive sentence says, "The city built public
transport and added green space." An argumentative sentence says, "Public
transport and green space can reduce congestion, emissions, heat stress, and
recreation inequality, but their success depends on affordability and whether
the network connects peripheral lower-income districts to jobs." The second
sentence is better because it contains outcome, process, stakeholder, and
limitation. That is the density you need in a \term{10-mark} style answer.

\textbf{Minute 2--4: social and environmental stresses.}
Urban stress means pressure on people, infrastructure, or the environment caused
by the way a city grows, functions, and is managed. Separate stresses before
you write. \term{Environmental stresses} include air pollution, congestion,
waste, urban heat island effects, flood risk, water pollution, loss of green
space, high energy use, and pressure on ecosystems at the urban fringe.
\term{Social stresses} include housing shortage, overcrowding, informal
settlement, segregation, inequality, displacement, service gaps, crime fears,
and unequal access to transport, healthcare, education, water, or sanitation.

The best answers do not list stresses. They explain why stresses appear.
Congestion may result from high car dependence, weak public transport, commuting
from suburbs, and road networks that cannot keep pace with growth. Housing
stress may result from rapid migration, high land values, limited formal
housing supply, poverty, and weak planning. Waste stress may result from high
consumption, insufficient collection, informal dumping, or weak recycling
systems. Flood stress may result from sealed surfaces, building on floodplains,
weak drainage, deforestation upstream, or climate-related heavier rainfall.

Now connect stress to inequality. A city can have world-class offices, airports,
and cultural districts while poorer households face long commutes, informal
housing, weak sanitation, or displacement. This matters because Urban
environments is not only about city form. It is about who gains access to the
benefits of urban life and who carries the costs. In an extended answer,
evaluation often begins with the question: \term{who benefits, who pays, and
who is left out?}

\textbf{Minute 4--6: sustainability as a response.}
Sustainable urban management tries to make cities more socially fair,
economically viable, and environmentally lower-impact over time. The three
classic pillars are social, economic, and environmental sustainability. In IB
answers, do not mention the pillars as decoration. Use them as criteria for
judgment.

For environmental management, possible strategies include public transport,
walking and cycling infrastructure, congestion charging, cleaner energy,
green buildings, recycling, waste reduction, flood-sensitive design, urban
trees, parks, wetlands, and protection of greenbelts or brownfield reuse. For
social management, strategies include affordable housing, slum upgrading,
tenant protection, public services, community participation, safer public
spaces, inclusive transport, and planning that connects lower-income districts
to jobs and services. For economic management, strategies include job creation,
regeneration, investment in infrastructure, support for informal livelihoods,
and redevelopment of derelict land.

The evaluation is that strategies can solve one problem while creating another.
A new metro line may reduce congestion and emissions, but if station-area land
values rise, poorer residents may be priced out. Waterfront regeneration may
improve the image of a city and attract investment, but it may also shift public
money toward tourists, developers, or higher-income residents. Greenbelts may
protect farmland and ecosystems, but they can also restrict housing supply and
increase prices if not managed with affordable housing policies. This is why a
strong answer judges trade-offs rather than calling a strategy simply
successful.

\textbf{Minute 6--8: future urban management debate.}
Future urban management is a debate because cities face linked pressures:
population growth, ageing infrastructure, housing costs, climate risk, resource
demand, inequality, and economic restructuring. There is no single solution.
Different visions of the future city emphasize different priorities.

Use the syllabus phrase \term{building sustainable urban systems for the
future} as the anchor. The city labels below are supplementary planning lenses,
not a separate IB syllabus list. They are useful only if they help you discuss
future sustainability, social inclusion, and environmental impact.

A \term{compact city} approach tries to limit sprawl by increasing density,
mixing land uses, improving public transport, and making services closer to
where people live. It can reduce car dependence and protect land at the edge,
but it may increase pressure on open space or housing costs if density is not
well planned. A \term{smart city} approach uses data, sensors, digital services,
and networked infrastructure to manage transport, energy, waste, and safety. It
can improve efficiency, but it raises questions about cost, access, privacy,
and whether poorer districts are included. A \term{resilient city} approach
focuses on preparing for shocks such as floods, heatwaves, disease outbreaks,
economic disruption, and infrastructure failure. It can protect vulnerable
groups, but resilience planning needs money, trust, and long-term governance.

The strongest future-management argument is usually conditional. For example:
"Sustainable urban management is most effective when environmental strategies
are paired with social inclusion." Or: "Technology can improve future cities,
but only where governance ensures that benefits reach lower-income residents."
Conditional judgments are better than one-sided claims because they show
geographical complexity.

Now add stakeholder conflict. A future urban strategy can look successful from
one viewpoint and limited from another. A commuter may value faster public
transport. A renter may fear rent increases near stations. A developer may
value higher land prices. A city government may value tax revenue and a stronger
image. A low-income household may value affordable access more than
architectural design. A climate planner may prioritize trees, drainage, and
lower emissions. An exam answer improves when you name these different users of
the city because sustainability is not only an environmental target; it is also
a question of fairness and power.

Add \term{ecological footprint} when the prompt asks about city sustainability.
A city's ecological footprint is the land, water, energy, materials, and waste
absorption needed to support its consumption. Dense cities can sometimes reduce
per-person transport and land consumption, but high-income urban lifestyles can
still have large footprints through electricity use, imported food, construction
materials, water demand, consumer goods, and waste. This gives you another
evaluation line: a strategy may make the visible city greener while leaving the
hidden resource footprint of consumption largely unchanged.

Use a judgment ladder. At the bottom is a weak claim: "The strategy is good."
One step up is: "The strategy reduces congestion." Better is: "The strategy
reduces congestion and emissions in connected areas." Stronger again is:
"The strategy reduces congestion and emissions in connected areas, but its
social success depends on affordability and whether lower-income residents can
use it." The top step is a conclusion that weighs the whole issue:
"Therefore the strategy is environmentally effective but only partly
sustainable unless social access is protected." Try to reach the top step in
your plan before you begin any full answer.

\textbf{Minute 8--10: examples and case-study discipline.}
Use your class examples first. The student case-study sheet lists
\term{Gentrification in Detroit}. Detroit is useful for deindustrialization,
population loss, vacancy, selective reinvestment, gentrification, regeneration,
and uneven benefits. A usable Detroit argument is: long-term industrial decline
and population loss created vacant land and social stress; later reinvestment
in selected districts such as Downtown or Midtown improved image, services, and
economic activity; however, benefits were uneven and some residents may face
rising rents, exclusion, or limited improvement in poorer neighborhoods.

Detroit is not the best example for every sustainability question. If the prompt
asks about sustainable transport, use the exact transport or planning case your
class has studied. If you have no class-specific case ready, use a placeholder
in the plan and mark it as a gap. A common practice anchor is \term{Curitiba,
Brazil}, often used for bus rapid transit, land-use planning, parks, and
recycling initiatives. If you use it, evaluate it: integrated planning can
reduce congestion and improve access, but success depends on funding, continued
maintenance, social inclusion, and whether peripheral low-income residents are
fully served. Do not memorize an example as a slogan. Store it as a usable
argument: problem, strategy, result, limitation.

The minimum evidence package for any urban essay is:
\term{place} $\rightarrow$ \term{problem} $\rightarrow$ \term{strategy or
process} $\rightarrow$ \term{stakeholders} $\rightarrow$ \term{impact}
$\rightarrow$ \term{limitation}. If one link is missing, write that gap in the
workbook before the May 1 case-study bank day.

Avoid three common mistakes. First, do not case-study dump. A long paragraph
about Detroit that never answers the prompt will not score as well as a shorter
Detroit paragraph tied directly to gentrification, inequality, and management
success. Second, do not treat sustainability as only "green". A strategy that
lowers emissions but displaces poorer residents is environmentally useful but
not fully sustainable. Third, do not end with a vague balanced sentence such as
"there are advantages and disadvantages." Replace it with a precise judgment:
"The strategy is more successful environmentally than socially because benefits
are concentrated in accessible, reinvested districts."

\textbf{Minute 10--12: model plan and judgment.}
Now build a model plan for this prompt: \exam{Evaluate the success of strategies
used to reduce social and environmental stresses in one or more urban areas.}
The introduction should define success using criteria. Success might mean
reduced pollution, improved transport, more affordable housing, lower flood
risk, better access to services, and fairer outcomes for lower-income
residents. The judgment could be: strategies can reduce selected stresses, but
they are only fully successful when environmental improvements are linked to
social equity and long-term governance.

Paragraph one could focus on environmental strategy. Public transport, green
space, waste management, or brownfield redevelopment can reduce congestion,
pollution, land consumption, or dereliction. Evaluation: these strategies may
work best in well-funded or central areas, and benefits may be spatially uneven.

Paragraph two could focus on social strategy. Affordable housing, service
investment, slum upgrading, or tenant protection can reduce overcrowding,
displacement, and inequality. Evaluation: these strategies are politically and
financially difficult because land values, private developers, and limited
public budgets can weaken inclusion.

Paragraph three should compare stakeholders or scales. A strategy may be
successful for city image, investors, commuters, or tourists, while being less
successful for informal workers, renters, elderly residents, or peripheral
households. This is where the answer becomes evaluative rather than descriptive.

The conclusion should not repeat the introduction. It should answer the exact
question: "Overall, urban management strategies are most successful when they
treat social and environmental stresses as linked. A transport or regeneration
project may improve the physical city, but if it raises land values or excludes
poorer residents, the success is partial rather than complete." That is the
kind of judgment to aim for in the 25-minute workbook block.

Pause now and say a five-part plan from memory:
\term{define success} $\rightarrow$ \term{environmental stress}
$\rightarrow$ \term{social stress} $\rightarrow$ \term{stakeholder evaluation}
$\rightarrow$ \term{clear final judgment}.

\subsection{Urban Essay Planning Tools}

\rowcolors{2}{TableStripe}{white}
\begin{tabularx}{\textwidth}{L{0.2\textwidth}Y Y}
\toprule
\textbf{Planning move} & \textbf{Question to ask} & \textbf{Useful answer language} \\
\midrule
Command term & Do I need to describe, explain, examine, evaluate, or judge
extent? & This answer will judge success by... \\
Criteria & What would count as success? & Success is partial unless it improves
both environmental quality and social access. \\
Case evidence & Which place proves the point? & In Detroit..., in the class
transport case..., in the named regeneration scheme... \\
Stakeholders & Who benefits and who may lose? & Benefits are uneven because
commuters, developers, renters, and informal workers experience the strategy
differently. \\
Scale & Is the strategy city-wide, district-level, or household-level? & The
strategy may work in the central district but not across the whole urban area. \\
Judgment & How far is the claim true? & Therefore the strategy is effective in
reducing one stress, but limited as a complete sustainability solution. \\
\bottomrule
\end{tabularx}
\rowcolors{2}{}{}

\subsection{Urban Stress And Management Matrix}

\rowcolors{2}{TableStripe}{white}
\begin{tabularx}{\textwidth}{L{0.2\textwidth}Y Y Y}
\toprule
\textbf{Stress} & \textbf{Cause} & \textbf{Possible response} & \textbf{Evaluation point} \\
\midrule
Congestion and air pollution & Car dependence, commuting, limited public
transport, decentralization & Buses, rail, cycling, walking, congestion charge,
mixed-use planning & Can reduce emissions, but low-income peripheral areas
need affordable access \\
Housing shortage and displacement & Rapid growth, high land values, weak
affordable supply, redevelopment pressure & Affordable housing, tenant
protection, service upgrading, brownfield housing & May be limited by cost,
private land markets, and political resistance \\
Waste and pollution & High consumption, weak collection, informal dumping,
industrial activity & Recycling, collection systems, circular economy,
regulation & Needs behavior change, funding, enforcement, and inclusion of
informal workers \\
Flood and heat risk & Sealed surfaces, low green space, floodplain building,
climate stress & Green infrastructure, parks, wetlands, shade, drainage, cool
roofs & Reduces risk, but benefits may concentrate in wealthier districts \\
Gentrification pressure & Reinvestment, cultural amenities, improved image,
limited housing protections & Inclusive zoning, rent support, community
planning, affordable units & Can improve place quality while still excluding
some original residents \\
\bottomrule
\end{tabularx}
\rowcolors{2}{}{}

\section{Audio 2 Script}

\subsection{Population Distribution -- Changing Population}

\textit{Target length: about 12 minutes at a natural revision pace. This is a
second-pass Population session, so focus on retrieval, comparison, and policy
judgment rather than passive rereading.}

\textbf{Minute 0--2: the four-word frame.}
The second April 29 audio returns to \term{Population distribution -- changing
population}. Use four words to organize the unit:
\term{pattern}, \term{process}, \term{structure}, and \term{policy}. Pattern
asks where people are and why they are unevenly distributed. Process asks how
population numbers change through births, deaths, and migration. Structure asks
what the age-sex composition looks like. Policy asks how governments try to
manage population change.

This frame is useful because Paper 2 questions often begin with a graph, map,
table, or population pyramid. If you panic, return to the four words. Describe
the pattern shown. Explain the process that may have created it. Infer the
structure or consequences. Then link to policy or planning. That sequence turns
unfamiliar data into a geographical answer.

Start with distribution. Population distribution is the spread of people across
space. Dense areas may be linked to flat land, fertile soils, reliable water,
coasts, rivers, trade routes, jobs, services, safety, and historical settlement.
Sparse areas may be linked to deserts, mountains, dense forests, cold climates,
poor access, hazards, conflict, or weak economic opportunity. But never explain
distribution with physical factors alone. Technology, wealth, trade, policy,
and infrastructure can make difficult environments more habitable, while
conflict or weak governance can limit settlement even in physically favorable
areas.

\textbf{Memory check.} Say the four-word frame. Then explain one dense
population pattern using one physical factor and one human factor.

Add scale to the frame. A pattern can be global, national, regional, urban, or
local. A country can have a moderate average population density while most
people are concentrated on a coast, river valley, plain, or capital region. A
city can grow overall while some inner districts lose population and outer
districts gain it. A rural region can lose young adults but retain older
residents. When you write, name the scale because population averages can hide
the variation that geography is trying to explain.

\textbf{Minute 2--5: change through births, deaths, and migration.}
Population change is built from a simple formula: \term{natural increase} plus
\term{net migration}. Natural increase means births minus deaths. Net migration
means immigration minus emigration. A population grows when births exceed deaths
and/or immigration exceeds emigration. A population declines when deaths exceed
births and/or emigration exceeds immigration. This formula is basic, but it
prevents vague answers.

\term{Fertility} is shaped by education, income, gender equality, urbanization,
child survival, access to contraception, religion, culture, housing cost, work
patterns, and state policy. Fertility often falls when child survival improves,
women gain more education and employment options, urban living becomes more
expensive, and families choose to have fewer children. However, the pace of
change varies because culture, inequality, healthcare, and policy vary.

\term{Mortality} is shaped by healthcare, nutrition, sanitation, conflict,
disease, disasters, age structure, and living conditions. Falling mortality
often reflects clean water, vaccines, better medical care, improved food supply,
and public health. Rising or high mortality can reflect conflict, epidemics,
poverty, environmental hazards, or ageing. If you write "better healthcare",
add the mechanism: it reduces deaths from disease, childbirth, or old-age
conditions, which raises life expectancy and changes population structure.

Migration links population to urban environments. Migrants may move because of
jobs, education, safety, family networks, services, conflict, environmental
stress, or inequality. Receiving areas may gain labor, skills, diversity, and
tax revenue, but may also face housing, service, and political pressure. Sending
areas may receive remittances, but they may lose young adults or skilled workers.
The evaluation is that migration produces gains and losses at different scales.

If your class uses the \term{demographic transition model}, use it carefully.
The model describes broad changes in birth rates, death rates, and total
population growth as societies develop, but it is not a perfect prediction for
every country. It helps explain why death rates may fall before birth rates,
creating rapid natural increase, and why later low fertility can lead to ageing.
Its limitation is that it can oversimplify culture, policy, migration, conflict,
inequality, and different development pathways. In an exam answer, the model is
useful only when it helps explain evidence; do not force every country into the
same sequence.

For stage-based interpretation, remember the broad pattern. Stage 1 has high
birth rates and high death rates, so growth is low and fluctuating. Stage 2 has
death rates falling faster than birth rates, so natural increase becomes rapid.
Stage 3 has falling birth rates, often linked to urbanization, education,
contraception, and changing costs of children. Stage 4 has low birth and death
rates, so growth is slow or stable. Stage 5 is sometimes used for places where
birth rates fall below death rates, creating natural decrease and ageing. Treat
these stages as a model for interpreting data, not as a rigid path every country
must follow.

\textbf{Minute 5--7: population structure and data interpretation.}
Population structure is the composition of a population by age and sex. The
main visual tool is the \term{population pyramid}. Read it in four steps:
\term{shape}, \term{process}, \term{consequence}, \term{policy}. A wide base
usually suggests high fertility and a youthful population. A narrow base with a
larger older section suggests low fertility, ageing, and longer life expectancy.
A working-age bulge can create a demographic dividend if education, health, jobs,
and governance allow people to work productively.

Do not treat a pyramid as a picture to describe and stop. Turn shape into
planning implications. A youthful structure increases demand for schools,
vaccination, maternal healthcare, housing, sanitation, and future jobs. An
ageing structure increases demand for pensions, healthcare, social care,
accessible transport, and labor-force planning. A migrant-heavy structure may
change housing demand, labor supply, language services, and political debate.

The key terms are \term{dependency ratio} and \term{demographic dividend}.
Dependency ratio compares dependent age groups with the working-age population.
A high youth dependency ratio points toward schools and future jobs. A high
old-age dependency ratio points toward pensions, healthcare, social care, and
tax pressure. Demographic dividend is the possible economic gain from a large
working-age share, but it is not automatic. It needs education, health, job
creation, gender inclusion, and stable governance.

Practise one data drill. Imagine Population A has a wide base, many children,
and very few elderly people. A good description is that Population A is
youthful. A good explanation is that fertility is likely high and life
expectancy may be lower than in an ageing society. A good consequence is that
schools, vaccination, maternal healthcare, water, sanitation, and future jobs
will be priorities. A good policy link is investment in education, health, and
employment, possibly with voluntary family-planning services. Now imagine
Population B has a narrow base and a large elderly share. Describe it as
ageing. Explain it through low fertility and high life expectancy. Consequences
include pensions, healthcare, social care, labor shortages, and possible rural
service decline. Policy may include later retirement, automation, immigration,
pronatalist support, or healthcare reform. That four-step method is how a visual
source becomes a complete answer.

\textbf{Minute 7--10: case-study comparison.}
The student case-study sheet gives three useful Population anchors:
\term{China's one-child policy}, \term{Japan's ageing population}, and
\term{Singapore's pronatalist and antinatalist policies}. Use them as a
comparison set.

China's one-child policy is an antinatalist example. It was designed to reduce
rapid population growth and ease pressure on food, housing, education, jobs, and
services. In a narrow sense, it helped reduce fertility and slow natural
increase. The evaluation is that it also contributed to ageing, a smaller future
workforce, gender imbalance, and ethical concerns about state control over
family size. The best sentence is balanced: it reduced one demographic pressure
while creating or intensifying other long-term social and economic pressures.

Japan's ageing population is useful for consequences of low fertility and long
life expectancy. An older age structure can increase pension costs, healthcare
demand, social care needs, labor shortages, and pressure on rural services. But
do not write as if older people are only a burden. Ageing can also create new
markets, community roles, demand for medical technology, and incentives for
automation or later retirement. The judgment is that ageing is a planning
challenge whose impact depends on productivity, migration policy, healthcare
systems, social attitudes, and government capacity.

Singapore is valuable because it shows population policy changing over time.
Policies can be antinatalist when rapid growth is seen as a problem and
pronatalist when low fertility and ageing become the problem. Pronatalist
measures may include financial incentives, childcare support, parental leave, or
housing support. The evaluation is that incentives do not automatically raise
fertility if housing costs, career pressure, gender expectations, delayed
marriage, or work culture remain strong. Singapore helps you write a more
sophisticated policy answer: policy operates within social and economic context.

Update the China policy timeline if the question asks how policy changed. The
one-child policy began nationally around 1979, was relaxed to a two-child policy
from 2015--2016, and shifted again to a three-child policy in 2021. The
geographical point is that a policy designed to slow rapid growth later had to
respond to low fertility, ageing, and labor-force concerns. This strengthens
evaluation because it shows that a population policy can be successful for one
period but less suitable when demographic structure changes.

Gender is not only a China add-on. It shapes fertility decisions, migration,
workforce participation, health, education, and vulnerability. Gender imbalance
can affect marriage patterns and social structure; gendered migration can change
who leaves rural areas or who works in care, domestic labor, construction, or
services; gender equality can reduce fertility by changing education, income,
and family-size choices. Use gender as a population variable whenever the source
or question points to age-sex structure, policy, migration, or household
decision-making.

Megacities and forced migration are two extra population topics to keep active.
A \term{megacity} has more than \term{10 million} people, and it matters because
very large urban populations intensify housing, transport, water, sanitation,
waste, employment, and governance challenges while also concentrating services,
innovation, and economic opportunity. \term{Forced migration} means movement
under pressure from conflict, persecution, disasters, development projects, or
environmental stress. It can create internally displaced people within one
country or refugees crossing borders. In answers, explain both origin impacts
and destination impacts: labor loss, remittances, housing pressure, service
pressure, social tension, humanitarian need, and possible cultural or economic
contribution.

Finally, connect population policy to \term{Malthusian} and
\term{Boserupian} thinking when useful. A Malthusian or neo-Malthusian view
emphasizes that population and consumption growth can outpace resources and
create pressure on food, water, jobs, housing, or the environment. A Boserupian
view argues that pressure can stimulate innovation, technology, and better
management. China's antinatalist policy can be discussed as a state response to
resource and development pressure, while Singapore's policy shifts show a more
adaptive view of population as both pressure and human capital.

Now compare the three examples directly. China shows a policy trying to reduce
growth. Japan shows the consequences of very low fertility and high longevity.
Singapore shows policy changing direction as the demographic problem changes.
Together they let you make a stronger argument than any one case alone:
population policy is not simply successful or unsuccessful; it is judged by
intended effects, unintended effects, ethics, and fit with wider society. That
comparison is especially useful for an \exam{evaluate} or \exam{examine}
question.

Also connect Population to the other core units. Rapid population growth can
increase demand for food, water, energy, housing, transport, and waste services.
Ageing can change consumption patterns, healthcare demand, tax systems, and
labor markets. Migration can be influenced by climate stress, conflict,
resource insecurity, or economic opportunity. This is why Population is not a
stand-alone unit. It is the human pressure and human capacity side of global
change.

\textbf{Minute 10--12: model Paper 2 answer habit.}
Here is a model paragraph for a population policy question. Population policies
are most effective when they address both the demographic target and the social
conditions behind it. China's one-child policy reduced fertility and slowed
rapid population growth, so it helped address pressure on services and
resources. However, the same policy contributed to long-term ageing, a smaller
future labor force, gender imbalance, and ethical concerns. Singapore's policy
history shows a different lesson: governments may reverse direction as
demographic priorities change, but pronatalist incentives are limited if high
living costs and work pressures still discourage larger families. Therefore,
population policy can influence demographic trends, but its success is partial
unless it fits wider social, economic, and cultural conditions.

Now practise the visual-stimulus habit that Paper 2 often needs. If a source
shows a map, begin with distribution: where are people concentrated, and where
are they sparse? Then explain likely physical and human reasons. If a source
shows a line graph, begin with trend: is the population, birth rate, death rate,
fertility rate, or migration rate rising or falling? Then suggest process. If a
source shows a population pyramid, begin with shape: base, middle, top, and sex
balance. Then infer fertility, mortality, life expectancy, migration, dependency,
and planning needs. The pattern should come first, but it should never be the
only thing you write.

For command terms, keep the response controlled. \exam{Describe} means say what
the source shows using clear comparative language. \exam{Explain} means give a
reason and show the mechanism. \exam{Suggest} means use source evidence and
plausible course knowledge. \exam{Compare} means write both similarity and
difference where possible. \exam{Evaluate} means judge success, limitation, or
importance. One useful sentence frame is: "The source shows..., this may be
caused by..., and this matters because..." That frame works for a population
map, graph, table, or pyramid.

Notice the structure. It compares examples, uses evaluation, and ends with a
judgment. That is the habit you need in Paper 2 and the same habit you need in
today's Paper 1 Urban extended-answer plan.

Finish with a retrieval drill. Say the formula for population change. Say the
four steps for reading a population pyramid. Name one youthful-population
consequence, one ageing-population consequence, one antinatalist policy
limitation, and one pronatalist policy limitation. Then link Population back to
Urban environments: population growth and migration can increase housing,
transport, water, sanitation, and service pressure; ageing can change transport,
healthcare, and housing needs; policy shapes who is supported and who is
excluded. That cross-unit thinking is what makes the core feel active rather
than memorized.

\subsection{Population Retrieval Table}

\rowcolors{2}{TableStripe}{white}
\begin{tabularx}{\textwidth}{L{0.23\textwidth}Y Y}
\toprule
\textbf{Term} & \textbf{Core meaning} & \textbf{Exam use} \\
\midrule
Population distribution & Spread of people across space & Explain with physical
and human factors together \\
Natural increase & Births minus deaths & Use to explain growth or decline
before adding migration \\
Net migration & Immigration minus emigration & Evaluate impacts on sending and
receiving areas \\
Population pyramid & Age-sex structure graph & Read shape, infer process, state
consequence, link to policy \\
Dependency ratio & Relationship between dependent age groups and working-age
population & Use for youthfulness, ageing, tax pressure, and service demand \\
Demographic dividend & Possible economic benefit from a large working-age
share & Evaluate using education, jobs, health, and governance \\
Antinatalist policy & Policy aiming to reduce fertility & Evaluate effects and
unintended consequences \\
Pronatalist policy & Policy aiming to raise fertility & Evaluate cost, culture,
housing, gender roles, and work pressure \\
Megacity & Urban area with more than 10 million people & Link population growth
to housing, transport, water, waste, employment, and governance pressure \\
Forced migration & Movement under pressure from conflict, persecution,
disaster, development, or environmental stress & Compare impacts on origin and
destination areas, including IDPs and refugees \\
Demographic transition model & Model of changing birth rates, death rates, and
growth through development stages & Use stages to interpret data, but evaluate
limits and avoid treating it as a fixed path \\
\bottomrule
\end{tabularx}
\rowcolors{2}{}{}

\subsection{Population Case Comparison}

\rowcolors{2}{TableStripe}{white}
\begin{tabularx}{\textwidth}{L{0.2\textwidth}Y Y}
\toprule
\textbf{Anchor} & \textbf{Use it for} & \textbf{Evaluation sentence} \\
\midrule
China's one-child policy & Antinatalist policy, fertility reduction, state
intervention & It reduced rapid growth, but contributed to ageing, gender
imbalance, and ethical concerns. \\
Japan's ageing population & Low fertility, long life expectancy, old-age
dependency & Ageing creates pension, healthcare, and labor challenges, but
outcomes depend on productivity, automation, migration, and social policy. \\
Singapore's changing policy & Antinatalist-to-pronatalist policy shift & Policy
can change direction, but incentives are limited if costs and work culture still
discourage births. \\
\bottomrule
\end{tabularx}
\rowcolors{2}{}{}

\section{25-Minute Practice Block}

Use the workbook file
\term{2026-04-29-urban-essay-planning-workbook.pdf}. The plan asks you to
choose one Paper 1 extended-answer prompt from a choice of two and produce a
clear plan with a judgment.

\subsection{Choice Of Two Prompts}

\boxnote{SoftTint}{Prompt A.}{\exam{Evaluate} the success of strategies used to
reduce social and environmental stresses in one or more urban areas.}

\vspace{0.5em}

\boxnote{SoftBlue}{Prompt B.}{\exam{To what extent} can sustainable urban
management solve the challenges of future urban growth?}

\subsection{The 25-Minute Rhythm}

\rowcolors{2}{TableStripe}{white}
\begin{tabularx}{\textwidth}{L{0.15\textwidth}Y Y}
\toprule
\textbf{Time} & \textbf{Task} & \textbf{Output} \\
\midrule
0--3 min & Closed-book recall of urban stresses, strategies, and examples &
Fast idea bank \\
3--6 min & Choose Prompt A or B and unpack command term & One focused question
translation \\
6--10 min & Decide judgment and criteria & One thesis sentence \\
10--18 min & Plan three main paragraphs & Evidence, explanation, evaluation
for each paragraph \\
18--22 min & Write conclusion in note form & Final judgment with degree of
success or extent \\
22--25 min & Self-check and fix weakest part & One improved judgment sentence \\
\bottomrule
\end{tabularx}
\rowcolors{2}{}{}

\subsection{End-Of-Day Output}

You are finished with April 29 if you have:

\begin{itemize}
  \item chosen one of the two Urban environments extended-answer prompts
  \item written a thesis that makes a clear judgment
  \item planned at least three paragraphs with evidence, explanation, and evaluation
  \item used Detroit or another named urban example accurately
  \item identified one missing urban management case detail to fill before May 1
  \item recalled population change using births, deaths, and migration
  \item compared China, Japan, and Singapore as population examples
  \item logged one improvement for the next essay-planning session on April 30
\end{itemize}

\boxnote{SoftGreen}{Carry forward.}{On Thursday 2026-04-30, the plan moves to
Food and health essay work plus Climate. Carry forward today's essay-planning
habit: do not write until the command term, criteria, evidence, and judgment are
clear.}

\end{document}
